\chapter{Antecedentes y Alcance}

\section{Contexto del Proyecto}

El sistema eléctrico colombiano presenta una alta dependencia de la generación hidráulica, la cual históricamente ha representado entre el 70\% y el 80\% de la energía eléctrica producida en el país. Esta concentración de la matriz energética expone al sistema a riesgos asociados a la variabilidad climática, como quedó evidenciado durante el fenómeno de El Niño de 2015-2016, periodo en el que los precios de la energía alcanzaron máximos históricos ante el riesgo de racionamiento.

Este escenario ha impulsado la necesidad de diversificar las fuentes de generación del país mediante tecnologías que no dependan de las condiciones hidrológicas. En línea con este propósito, el Estado colombiano ha promovido el desarrollo de fuentes no convencionales de energía renovable (FNCER) a través de instrumentos regulatorios como:

\begin{itemize}
    \item Ley 1715 de 2014
    \item Resolución CREG 060 de 2019
\end{itemize}

\section{Alcance del Estudio}

En este contexto, el supermercado MAS X MENOS sede Guarín adelanta el desarrollo del proyecto SOLAR MAS X MENOS GUARÍN, de 141,6 kWp de capacidad instalada, cuya conexión al Sistema de Distribución Local (SDL) requiere la elaboración de un estudio de conexión simplificado, conforme a lo dispuesto en la Resolución CREG 030 de 2018.

El presente documento desarrolla dicho estudio a partir de los requerimientos técnicos y la información de red suministrados por el operador de red (OR) correspondiente al punto de conexión del proyecto, en este caso ESSA.

\section{Objetivo General}

El objetivo general del presente estudio es analizar la conexión para el desarrollo del proyecto fotovoltaico de 141,6 kWp de capacidad en el nodo 3272966 de ESSA a un nivel de tensión de 13,8 kV. Para ello, se realizaron análisis técnicos que determinen los impactos de la nueva generación.

\section{Objetivos Específicos}

\begin{enumerate}
    \item Evaluar el desempeño del sistema eléctrico ante la conexión de la nueva generación.
    \item Determinar posibles sobrecargas en líneas y transformadores.
    \item Verificar que no haya violaciones de los límites de tensión establecidos.
    \item Analizar las corrientes de cortocircuito trifásico y monofásico.
    \item Validar la coordinación de protecciones en el punto de conexión.
    \item Calcular las pérdidas técnicas del sistema.
\end{enumerate}
